\section{Properties, thermodynamic provenance and units}
\label{app:properties}
This section lists the helium properties, molar masses, thermodynamic
records and optional PbI$_2$--He diffusivity used in the calculations.

\subsection{Helium and transported molecular masses}
The executed helium carrier uses
$M_{\rm He}=0.0040026\units{kg\,mol^{-1}}$,
$C_p=5193\units{J\,kg^{-1}\,K^{-1}}$ and
\begin{equation}
 \begin{gathered}
 \rho=\frac{pM_{\rm He}}{RT},\qquad
 \mu(T)=\frac{A_s\sqrt T}{1+T_s/T},\\
 A_s=1.4683\times10^{-6},\qquad T_s=83.40\units{K}.
 \end{gathered}
\end{equation}
The numerical value of $A_s$ gives viscosity in Pa\,s when $T$ is in K.
These are case inputs to the perfect-gas, constant-heat-capacity,
Sutherland model, not helium properties newly fitted to measurements.
Partial-pressure conversions use absolute pressure, while
thermodynamic activities use $p^\circ=10^5\units{Pa}$. The assumed
controller reference pressure, 101,325 Pa, is a separate quantity.
\Cref{tab:masses} lists the species masses used in transport and
stoichiometric conversion.
\begin{table}[htbp]
\centering
\caption{Molar masses of the nine transported gases. Condensed iodide
inventories use the same formula masses; the shared Pb and Bi
inventories use atomic masses.}
\label{tab:masses}
\begin{tabular}{lr@{\qquad}lr@{\qquad}lr}
\toprule
Gas & kg/mol & Gas & kg/mol & Gas & kg/mol \\
\midrule
Pb & 0.207200 & Bi & 0.2089804 & I & 0.12690447 \\
PbI & 0.33410447 & Bi$_2$ & 0.4179608 & I$_2$ & 0.25380894 \\
PbI$_2$ & 0.46100894 & BiI & 0.33588487 & BiI$_3$ & 0.58969381 \\
\bottomrule
\end{tabular}
\end{table}

\subsection{Thermodynamic records and phase treatment}
Pb--I--He functions are evaluated from external NASA nine-coefficient
records labelled with Gurvich source information. Bi--I functions use
an external ThermoFun compilation labelled HERACLES/Barin. A common
property evaluator converts all records to a consistent absolute
convention, $G^\circ(T)=H^\circ(T)-TS^\circ(T)$, before forming
reaction Gibbs energies. Mixing this convention with an apparent
formation Gibbs energy would give inconsistent equilibrium constants.
The solver tables store formation constants and phase-resolved vapour
pressures as logarithms. Input-file hashes, record identifiers and table-generation
scripts are recorded in the benchmark provenance.

For an iodide with one gas molecule per condensed formula unit,
\begin{equation}
 p_{\rm eq}(T)=p^\circ\exp\left[-\frac{G_g^\circ(T)
                                  -G_c^\circ(T)}{RT}\right],
 \label{eq:pv}
\end{equation}
evaluated on the stable available condensed branch. The BiI$_3$ liquid branch
is constructed at $T_{\rm fus}=680.85\units{K}$ from the crystalline
record with $\Delta H_{\rm fus}=31.798\units{kJ\,mol^{-1}}$ and
$C_{p,l}=150.624\units{J\,mol^{-1}\,K^{-1}}$, as recorded in the
input-generation metadata. It is a constructed phase record, not an
independently confirmed liquid thermodynamic dataset. Bi and
Bi$_2$ share condensed Bi but have different gas chemical potentials
and stoichiometry.

The BiI $+29$ kJ/mol variant modifies the gas record and regenerates
all formation and disproportionation constants consistently. Both the
baseline and the perturbed functions are hypotheses for the
sensitivity calculation; the BiI source record still needs independent
confirmation. The external coefficient files are not reproduced in the
manuscript data package. Numerical conservation of the model derived
from them cannot establish their physical correctness or complete the
provenance assessment.

\subsection{Optional \texorpdfstring{PbI$_2$}{PbI2}--He diffusion estimate}
Only the PbI$_2$--He transport variant uses the following implemented
correlation:
\begin{align}
 D_{\rm PbI_2-He}&=
 \frac{1.86\times10^{-7}T^{3/2}\sqrt{1/4.0026+1/461}}
 {(p/101325)\sigma^2\Omega(T^*)},\\
 T^*&=T/71.117857,\qquad \sigma=4.258468,\\
 \Omega(T^*)&=0.5313-1.2946e^{-0.9922T^*}\nonumber\\
             &\quad+\frac{1.5591}{T^*}-\frac{0.1918}{(T^*)^2}.
\end{align}
In this numerical expression $T$ is in K, $p$ in Pa, $\sigma$ in
\AA\ and the masses in g/mol; the resulting $D$ is in
$\mathrm{m^2\,s^{-1}}$. To reproduce the original correlation, the
implementation keeps 461 g/mol rather than the more precise transport
formula mass in \cref{tab:masses}. Its estimated collision parameters
follow the original polarisability/electron-count construction
recorded in the code. The other eight gases keep the constant baseline
coefficient in this variant. The correlation is a screening alternative,
not an experimentally validated binary diffusivity or a replacement
for all missing Pb--Bi--I transport data.
